\section{The \texttt{MaybeUncertain<T>} Type}
\label{sec:type}

This section defines the type at the abstract level: its mathematical
structure, the four constructors, the sampling semantics, the operator
algebra, and the gating operator. Rust API details and the surrounding
machinery appear in~\Cref{sec:implementation}.

\subsection{Definition}
\label{sec:type:def}

A \texttt{MaybeUncertain<T>} value is an element of the product
\begin{equation}
\Mtype\langle T \rangle \;=\; \Utype\langle \Bool \rangle \,\times\, \Utype\langle T \rangle,
\label{eq:mtype}
\end{equation}
where $\Utype\langle T \rangle$ denotes the Bornholt-style uncertain type
over base type $T$. We write a value as a pair $(x, v)$, with $x \in
\Utype\langle \Bool \rangle$ the \emph{presence} component and $v \in
\Utype\langle T \rangle$ the \emph{value} component. The presence component
is a Bernoulli-distributed random variable governing whether the value
exists; the value component carries the distribution of the value
conditional on its existence.

\paragraph{Independence assumption.} The two components of a
$\Mtype\langle T \rangle$ value are independent random variables. In Rubin's
terms~\cite{rubin1976missing}, this restricts the semantics to data that is
missing at random: the event of presence carries no information about the
value the variable would take if observed. We state the assumption
explicitly here and revisit its consequences for missing-not-at-random
workloads in~\Cref{sec:discussion}.

\paragraph{Specialisations.} The library realises two instantiations,
$\Mtype\langle f64 \rangle$ and $\Mtype\langle \Bool \rangle$. The former
carries the full operator algebra described in~\Cref{sec:type:algebra};
the latter provides constructors, sampling, the presence predicates, and the
gating operator only. We note this asymmetry as a scope limitation
in~\Cref{sec:discussion}, not as a property of the type theory.

\subsection{Constructors}
\label{sec:type:ctors}

Four constructors cover the situations that arise in sparse-data programs.
Each constructor specifies the presence component and the value component
independently. Let $\delta_b$ denote the point distribution at $b$, and let
$\mathrm{Bern}(p)$ denote the Bernoulli distribution with success
probability $p$.

\begin{table}[h]
\centering
\small
\begin{tabular}{@{}lcc@{}}
\toprule
Constructor & Presence component $x$ & Value component $v$ \\
\midrule
$\textsc{fromValue}(c)$              & $\delta_{\text{true}}$ & $\delta_c$ \\
$\textsc{fromUncertain}(u)$          & $\delta_{\text{true}}$ & $u$ \\
$\textsc{alwaysNone}()$              & $\delta_{\text{false}}$ & $\bot$ \\
$\textsc{fromBernoulli}(p,\, u)$     & $\mathrm{Bern}(p)$ & $u$ \\
\bottomrule
\end{tabular}
\caption{The four constructors of $\Mtype\langle T \rangle$ and the
distributions they assign to each component. The placeholder $\bot$ in
$\textsc{alwaysNone}$ is never observed because the sampling semantics
short-circuit when the presence draw is false (\Cref{sec:type:sampling}).}
\label{tab:constructors}
\end{table}

The four constructors model four distinct situations. \textsc{fromValue}
represents a fact: the value exists with certainty and equals a known
scalar. \textsc{fromUncertain} represents a noisy reading that was definitely
received; presence is certain, value is uncertain. \textsc{alwaysNone}
represents structural absence. \textsc{fromBernoulli} is the canonical
sparse-data case: presence is governed by a Bernoulli trial, and the
caller supplies the conditional value distribution.

\subsection{Sampling semantics}
\label{sec:type:sampling}

Sampling a $\Mtype\langle T \rangle$ value $(x, v)$ produces an element of
$\{\bot\} \cup T$ by the two-stage procedure shown
in~\Cref{alg:sample}. The presence component is sampled first; the value
component is sampled only when the presence draw yields true.

\begin{figure}[h]
\centering
\fbox{\parbox{0.65\linewidth}{%
\textbf{Algorithm: $\mathrm{sample}(x, v)$.} \\[2pt]
\textbf{Input:} a $\Mtype\langle T \rangle$ value $(x, v)$. \\
\textbf{Output:} an element of $\{\bot\} \cup T$. \\[4pt]
\hspace*{1em} 1.\; Draw $b \sim x$ from the presence component. \\
\hspace*{1em} 2.\; \textbf{if} $b = \text{true}$ \textbf{then} draw $t \sim v$ and \textbf{return} $t$. \\
\hspace*{1em} 3.\; \textbf{else return} $\bot$.
}}
\caption{Two-stage sampling. The value component is sampled only on the
presence-true branch.}
\label{alg:sample}
\end{figure}

The induced output distribution is a mixture with an atom at $\bot$. Let
$p = \Prob[x = \text{true}]$. Then for any measurable $S \subseteq T$,
\begin{equation}
\Prob[\mathrm{sample}(x, v) \in S] \;=\; p \cdot \Prob[v \in S \mid x = \text{true}],
\quad \Prob[\mathrm{sample}(x, v) = \bot] \;=\; 1 - p.
\label{eq:mixture}
\end{equation}
Under the independence assumption, the conditional $\Prob[v \in S \mid x =
\text{true}]$ equals $\Prob[v \in S]$, so~\eqref{eq:mixture} simplifies to
$p \cdot \Prob[v \in S]$ on the value side. Repeated sampling treats each
draw as an independent realisation of the joint random variable;
deterministic re-sampling for tests and for the SPRT-gated operator is
discussed in~\Cref{sec:semantics}.

\subsection{Operator algebra}
\label{sec:type:algebra}

The $f64$ specialisation supports five operators: $+$, $-$, $\times$,
$\div$, and unary negation. Every binary operator follows the same pattern.
For an operator $\oplus$ acting on $\Utype\langle f64 \rangle$, the lifted
operator on $\Mtype\langle f64 \rangle$ is defined by
\begin{equation}
(x_1, v_1) \oplus (x_2, v_2) \;=\; \bigl(\,x_1 \wedge x_2,\; v_1 \oplus v_2\,\bigr),
\label{eq:algebra}
\end{equation}
where $\wedge$ denotes the logical-AND operator on
$\Utype\langle \Bool \rangle$. Unary negation acts on the value component
only and leaves the presence component unchanged:
\begin{equation}
-(x, v) \;=\; (\,x,\; -v\,).
\label{eq:neg}
\end{equation}

\paragraph{Choice of AND on the presence axis.} The result of a binary
operation between two $\Mtype\langle T \rangle$ values exists only if both
operands existed. We considered the dual choice in which absent operands
act as a neutral element under OR; that rule produces values whose
existence is invented by the operation itself, and contradicts the basic
sanity check that adding to an absent value should not summon one into
existence. We adopt~\eqref{eq:algebra} on those grounds.

\paragraph{Boolean specialisation.} The $\Bool$ specialisation does not
currently provide a lifted operator algebra. Compound presence-aware
Boolean expressions require operating on the presence and value components
separately; we note the gap in~\Cref{sec:discussion}.

\subsection{Predicates and the gating operator}
\label{sec:type:predicates}

Three additional operations connect $\Mtype\langle T \rangle$ to the
existing $\Utype$ infrastructure.

\paragraph{Presence predicates.} The operators $\textsc{isSome}$ and
$\textsc{isNone}$ project a $\Mtype\langle T \rangle$ value onto its
presence component and its negation, respectively:
\begin{equation*}
\textsc{isSome}(x, v) = x, \qquad \textsc{isNone}(x, v) = \neg x.
\end{equation*}
Both return $\Utype\langle \Bool \rangle$ values, recoverable by the existing
machinery; in particular, callers can compose them with the standard
Boolean operators on $\Utype\langle \Bool \rangle$ to express
presence-aware conditions over multiple $\Mtype\langle T \rangle$ values.

\paragraph{Lifting to $\Utype\langle T \rangle$.} The gating operator
$\textsc{lift}$ converts a $\Mtype\langle T \rangle$ value into a plain
$\Utype\langle T \rangle$ value, but only when statistical evidence
supports the presence hypothesis. Its signature is
\begin{equation*}
\textsc{lift} : \Mtype\langle T \rangle \,\times\, [0,1] \,\times\, [0,1] \,\times\, [0,1] \,\times\, \mathbb{N} \;\to\; \Utype\langle T \rangle \,\cup\, \{\mathrm{PresenceError}\},
\end{equation*}
with arguments $(x, v)$, threshold $p_\tau$, confidence $1 - \alpha$,
indifference half-width $\varepsilon$, and sample cap $N$. The operator
runs a sequential probability ratio test on the presence component $x$
against the hypotheses
\begin{equation*}
H_0 : \Prob[x = \text{true}] \le p_\tau - \varepsilon
\qquad\text{versus}\qquad
H_1 : \Prob[x = \text{true}] > p_\tau + \varepsilon,
\end{equation*}
at Type I and Type II error budgets each equal to $\alpha$. If the test
accepts $H_1$, the operator returns the value component $v$; if it accepts
$H_0$, it returns a $\mathrm{PresenceError}$. The full SPRT machinery is
described in~\Cref{sec:semantics:sprt}.

\paragraph{Strict and lax gating.} The same operator supports very different
operating points by choice of parameters. A high threshold with high
confidence and a large sample cap yields strict gating for safety-critical
paths; a low threshold with modest confidence yields lax gating suitable
for exploratory analysis. The decision is at the call site, not at the
type. We exploit this flexibility in the worked examples
(\Cref{sec:examples}).
